% =============================================================================
% Chapter 5: Architecture
% =============================================================================

\chapter{Architecture}
\label{ch:architecture}

\section{Module Structure}

\itc{} is organized into eight top-level packages, each with a single,
well-defined responsibility. The dependency graph is strictly layered:
\code{core} has no intra-project dependencies; everything else builds on
\code{core}. The \code{aerospace} package builds on \code{pproc}, which
builds on \code{ops}, \code{io}, and \code{uncertainty}.

\section{Full Module Tree}

\begin{lstlisting}[language=]
itaca/
+-- core/
|   +-- varframe.py            # VarFrame -- central structure
|   +-- uncframe.py            # UncFrame -- uncertainty mirror
|   +-- historyframe.py        # HistoryFrame -- origin tags (0/+1/-1)
|   +-- dimension.py           # Dimension dataclass with optional unit
|   +-- variable.py            # Variable dataclass with optional unit
|   +-- correlation.py         # CorrelationMatrix declaration & storage
|   +-- index.py               # dimension and index management
|   +-- coords.py              # Cartesian and Polar spatial coordinates
|   +-- axes.py                # Axis, AxisRegistry, rotate, translate_moments
|   +-- provenance.py          # Provenance: where data came from
|   +-- history.py             # History: what happened to the data
|   +-- pipeline.py            # Pipeline: serializable history subrange
|   +-- combine.py             # db.combine() arithmetic combination
|   +-- arithmetic.py          # internal arithmetic with covariance
|   +-- errors.py              # ITACAError hierarchy
|   `-- version.py             # version tagging
+-- io/
|   +-- loader.py              # itc.load() -- all input modes
|   +-- inspector.py           # db.inspect() on any datapoint-mode VarFrame
|   +-- pivot.py               # db.pivot() on any datapoint-mode VarFrame
|   +-- manifest.py            # db.manifest()
|   +-- diagnostics.py         # db.diagnostics() -> DiagnosticsReport
|   +-- summary.py             # db.summary() -- short, practical
|   +-- export.py              # to_csv, to_json, to_pandas, to_numpy, save
|   +-- prov_export.py         # PROV-N and PROV-JSON serialization
|   `-- formats/
|       +-- csv.py
|       +-- json.py
|       +-- pandas_bridge.py
|       +-- numpy_bridge.py
|       `-- itc.py             # native binary .itc format
+-- ops/
|   +-- select.py              # select() with operators and Frame= keyword
|   +-- squeeze.py             # squeeze() -- remove unit-cardinality dims
|   +-- expand.py              # expand() -- ADD a new dimension
|   +-- concat.py              # itc.concat() -- join VarFrames along a dim
|   +-- interpolate.py         # interpolate() -- densify or translate axis
|   +-- fill.py                # fill() -- gap fill via window or global
|   +-- compute.py             # compute() -- string equation derivation (RPN)
|   +-- average.py             # average() -- collapse dim by mean
|   +-- integrate.py           # integrate() -- Cartesian and polar
|   +-- smooth.py              # smooth() -- savgol, spline, moving_avg
|   +-- diff.py                # diff() -- moving polynomial derivative
|   +-- fitmodel.py            # fitmodel() -- replace dim with poly coefs
|   `-- fitvalue.py            # fitvalue() -- evaluate fit at new x
+-- uncertainty/
|   +-- expression.py          # ast-based parser, chain-rule differentiation
|   +-- operators.py           # operators with partial derivatives
|   +-- propagation.py         # GUM clause 5 LPU with covariance
|   `-- mcm.py                 # Monte Carlo for discrete-branch models
+-- pproc/
|   +-- factory.py             # itc.pproc() -- factory function
|   +-- protocol.py            # Processor typing.Protocol
|   +-- base.py                # Processor base class
|   +-- registry.py            # built-in processor registry
|   +-- statistics.py          # pproc.statistics()
|   +-- compare.py             # pproc.compare() -- N VarFrames, named inputs
|   +-- report.py              # pproc.report() -- LaTeX-backed PDF
|   +-- builtin/
|   |   +-- cfd_cp.py
|   |   +-- cfd_disk.py
|   |   +-- wt_balance_off.py
|   |   +-- wt_balance_on.py
|   |   +-- wt_propeller.py
|   |   +-- wt_corrections.py  # solid_blockage, streamline, buoyancy, ...
|   |   +-- ft_thrust.py       # SAE AIR 4065
|   |   +-- cal_balance.py
|   |   `-- aero_polar.py
|   `-- equations/
|       +-- parser.py          # .itceq parser with cycle detection
|       +-- balance/
|       +-- propeller/
|       +-- flight_test/
|       `-- cfd/
+-- surrogate/
|   +-- smt_bridge.py          # interface with SMT
|   `-- models.py              # itc.surrogate() returning (VarFrame, Surrogate)
+-- aerospace/
|   +-- performance/
|   |   +-- takeoff.py
|   |   +-- climb.py
|   |   `-- envelope.py
|   +-- stabcontrol/
|   |   +-- static.py          # static margin, trim + EUB
|   |   +-- damping.py         # pitch damping derivatives
|   |   `-- vn_diagram.py
|   `-- propeller/
|       +-- bemt.py            # Adkins-Liebeck BEMT
|       `-- integration.py
+-- plot/
|   +-- itcfigure.py           # ItcFigure wrapper
|   +-- datavis.py             # DataVis class -- itc.datavis()
|   +-- multiplot.py           # multi-figure layout
|   +-- backends/
|   |   +-- matplotlib_backend.py  # default backend
|   |   `-- plotly_backend.py      # optional, interactive
|   +-- style.py               # AIAA rcParams
|   +-- cycles.py              # color, marker, linestyle cycles
|   `-- specialized/
|       +-- polar.py
|       +-- cp.py
|       +-- disk.py
|       `-- colormap.py
`-- utils/
    +-- units.py               # custom unit-conversion table
    +-- memory.py              # sparse-aware helpers, footprint reporting
    `-- validation.py          # input validation and error formatting
\end{lstlisting}

\section{Top-Level API Surface}

Table~\ref{tab:toplevel-api} enumerates every function exposed at the
\code{itc.} top level, providing a single audit point for the public API.

\begin{table}[H]
\centering
\caption{Top-level \itc{} API surface.}
\label{tab:toplevel-api}
\small
\begin{tabularx}{\textwidth}{@{}lX@{}}
\toprule
\textbf{Function} & \textbf{Purpose}\\
\midrule
\code{itc.load(...)}            & Unified loader: folder, dict, NumPy array, pandas DataFrame, single file.\\
\code{itc.open(path)}           & Read a \code{.itc} archive into a VarFrame.\\
\code{itc.concat(list, along=)} & Concatenate compatible VarFrames along a dimension.\\
\code{itc.pproc(name\_or\_path)} & Construct a processor by registered name or \code{.itceq} path.\\
\code{itc.datavis(...)}         & Construct a DataVis plotting object.\\
\code{itc.surrogate(...)}       & Fit a surrogate, returning \code{(VarFrame, Surrogate)}.\\
\code{itc.load\_pipeline(path)} & Read a \code{.itc\_pipe} file.\\
\code{itc.load\_surrogate(path)}& Read a \code{.itc\_surr} file.\\
\code{itc.set\_user(name)}      & Override the default \code{user@hostname}.\\
\code{itc.set\_mode(mode)}      & Set the global default operating mode.\\
\code{itc.aerospace.*}          & Aerospace analysis subpackage (performance, stabcontrol, propeller).\\
\bottomrule
\end{tabularx}
\end{table}

\section{Class Hierarchy}

\itc{}'s principal classes follow a small, deliberately flat hierarchy.
The top-level data classes are \code{VarFrame}, \code{UncFrame}, and
\code{HistoryFrame}; each is a frozen composition of NumPy arrays plus
metadata classes (\code{Provenance}, \code{History}, \code{AxisRegistry},
\code{CorrelationMatrix}). Spatial coordinates split into \code{Cartesian}
and \code{Polar} via the \code{CoordSystem} abstract base. \code{Axis} is a
standalone dataclass usable across multiple VarFrames.

The processing layer is built on the \code{Processor} typing protocol; each
built-in processor (\code{WTPropeller}, \code{CFDDisk}, \code{AeroPolar},
\code{ItceqProcessor}, etc.) implements the protocol without subclassing a
common base. RPN operators (\code{Arith}, \code{Trig}, \code{Misc}) are
isolated, individually testable, and expose \code{evaluate}, \code{d\_da},
and \code{d\_db} methods for chain-rule differentiation.

\code{itc.aerospace} analyses (\code{TakeoffAnalysis},
\code{StabilityAnalysis}, \code{BEMTAnalysis}) are top-level callable
classes. Visualization is rooted at \code{DataVis}, which produces
\code{ItcFigure} objects. Surrogate models are wrapped in a single
\code{Surrogate} class with multiple backends.

\section{Key Design Decisions}

\begin{ddbox}{DD-02}{NumPy-only core}
The \code{core/}, \code{ops/}, and \code{uncertainty/} packages depend
exclusively on NumPy. No xarray, dask, or pandas imports appear in these
modules. All other integrations are optional and loaded lazily.
\end{ddbox}

\begin{ddbox}{DD-03}{VarFrame structural immutability}
\code{VarFrame} is a frozen dataclass. Operations return new objects.
In-place mutation is structurally impossible. This enables safe chaining
and ensures Provenance / History consistency.
\end{ddbox}

\begin{ddbox}{DD-04}{Processors as a Protocol, not a base class}
\code{Processor} is a \code{typing.Protocol} (structural subtyping).
Specialized analyses are callable objects called as \code{pproc(db)},
keeping data structure and analysis logic separate. Any object satisfying
the protocol is a valid processor; subclassing \code{Processor} is a
convenience, not a requirement.
\end{ddbox}

\begin{ddbox}{DD-05}{UncFrame as structural mirror}
Uncertainty is stored in a separate \code{UncFrame}, not as extra columns.
This keeps primary data clean and makes uncertainty presence explicit
(\code{db.uncertainty} is \code{None} when no uncertainties are assigned).
\end{ddbox}

\begin{ddbox}{DD-06}{HistoryFrame as separate mirror}
Origin tags ($0$/$+1$/$-1$) are stored in a separate \code{HistoryFrame},
not as extra columns or sentinel values within the main arrays. This keeps
the numerical arrays clean and allows operations to manipulate values and
tags in parallel without coupling.
\end{ddbox}

\begin{ddbox}{DD-07}{Operating modes as metadata, not API branch}
\code{draft} vs.\ \code{production} is a VarFrame attribute, not a
different class. The same API applies in both modes; only history-recording
behavior and export guards differ. Mixing modes in a binary operation is
forbidden and raises \code{OperatingModeMixError}; the user must call
\code{db.promote(mode="production")} or \code{db.demote(mode="draft")}
explicitly.
\end{ddbox}

\begin{ddbox}{DD-08}{Aerospace generates, pproc processes}
\code{itc.aerospace} modules \emph{generate} VarFrame data from physical
models (BEMT, performance equations, stability buildups). \code{pproc}
modules \emph{process} existing VarFrame data. This distinction is
fundamental to the architecture and enables a single suite of plotting,
comparison, and uncertainty machinery to apply uniformly to both.
\end{ddbox}

\begin{ddbox}{DD-09}{Sensitivity analysis as native capability}
Sensitivity analysis is not a module: it is the direct output of
combining \code{set\_uncertainty} and (optionally)
\code{set\_correlation} on input variables with \code{compute} or
\code{itc.aerospace} computations. The propagated uncertainty on the
output is the sensitivity measure.
\end{ddbox}

\begin{ddbox}{DD-10}{Errors as a typed hierarchy}
All \itc{}-specific exceptions inherit from a base class
\code{ITACAError}, organized into families: \code{DataError},
\code{ProcessorError}, \code{ProvenanceError}, \code{UncertaintyError},
\code{DependencyError}, \code{AxesError}. Users can catch family-level
exceptions when appropriate; specific subclasses provide actionable error
messages. The complete enumeration is given in
Table~\ref{tab:exception-hierarchy}.
\end{ddbox}

\begin{ddbox}{DD-11}{Custom axes are core, not aerospace}
\code{Axis} and the rotation operations live in \code{core/axes.py}, not in
\code{itc.aerospace}. Axis transformations are common across
experimental, CFD, and engineering workflows (WT bookkeeping routinely moves data between rig, tunnel, body, stability, and wind axes) and the
machinery must remain available outside the aerospace subpackage. Rotation
also interacts directly with the uncertainty propagation engine and
therefore lives next to \code{core/correlation.py}.
\end{ddbox}

\begin{ddbox}{DD-12}{Combine over operator overloading}
Arithmetic combinations of VarFrames go through
\code{db.combine(other, op=...)}, not Python operators. This forces every
combination to declare its operation explicitly, records the call in
History, and consults the declared correlation structure. The convenience
of \code{db1 + db2} is rejected as misleading (NREQ-08).
\end{ddbox}

\begin{ddbox}{DD-13}{Custom unit conversion, no external dependency}
Unit metadata on \code{Dimension} and \code{Variable} is optional and
opt-in. Conversion is implemented in \code{utils/units.py} with a
hand-curated table of SI and aerospace-engineering units. Pulling in
\code{pint} or \code{astropy.units} is rejected for v0.x to keep the
dependency surface minimal and the conversion table fully auditable.
\end{ddbox}

\begin{ddbox}{DD-14}{Covariance support from the first release}
The propagation engine always evaluates the full GUM clause-5 formula,
which reduces to the independence form when the correlation matrix is
zero. Wind tunnel multi-channel calibrations produce correlated
uncertainties as the norm, not the exception; deferring covariance would
ship an engine that systematically underestimates combined uncertainty for
the dominant use case. The correlation matrix is materialized lazily, so
users who declare no correlations pay no cost.
\end{ddbox}

\begin{ddbox}{DD-15}{Monte Carlo restricted to discrete-branch models}
\code{method="mcm"} exists specifically for expressions containing
discrete branches, where the linearization behind symbolic propagation
breaks down. Continuous nonlinearities are handled correctly by the
symbolic chain rule; switching them to Monte Carlo would be slower without
being more accurate.
\end{ddbox}

\begin{ddbox}{DD-16}{Idempotence opt-in, with explicit override}
Processors declare \code{idempotent: bool}, default \code{False}: a second
application warns and refuses unless \code{force=True} is passed, and a
permitted reapplication is always recorded as a distinct History entry.
Silent no-ops hide bugs; silent re-runs corrupt data (corrections applied
twice); warn-and-refuse protects against both.
\end{ddbox}

\begin{ddbox}{DD-17}{\code{.itceq} evaluation order is explicit}
Equations evaluate in file order by default; topological sorting is
opt-in via \code{auto\_sort=True} and reports the resolved order as
feedback. Forcing topological sort by default would make file behavior
depend on parser tiebreaking rules, a portability risk.
\end{ddbox}

\begin{ddbox}{DD-18}{Uncertainty semantics are defined per operation}
Every operation declares its effect on the UncFrame; the normative table
is REQ-98. All v0.x operations are linear in the variable values
(selection, interpolation, averaging, integration, smoothing kernels,
moving polynomial fits, rotation), so propagation through the operation
weights is exact. An operation with no sound propagation rule must raise
or warn; silently dropping or silently carrying the UncFrame are both
forbidden. The rejected alternative, propagation only in \code{compute},
left every other operation free to corrupt the uncertainty record.
\end{ddbox}

\begin{ddbox}{DD-19}{Two-component uncertainty}
The UncFrame separates systematic and random components
(Section~\ref{sec:two-components}). A single per-variable scalar was
rejected because reductions cannot then distinguish what shrinks with
$1/\sqrt{N}$ from what does not: repeat averaging would silently shrink
calibration bias, understating the combined uncertainty of every averaged
result, exactly the failure AIAA S-071A-1999 bookkeeping exists to
prevent.
\end{ddbox}

\begin{ddbox}{DD-20}{Expression parsing built on the Python ast module}
\code{db.compute} expressions are parsed with the standard-library
\code{ast} module and compiled to the \itc{} operator tree; symbolic
differentiation walks that tree. A hand-written tokenizer and parser was
rejected: it would be the most defect-prone component of the foundation
release, and \code{ast} is battle-tested, stdlib-only (the NumPy-only rule
is untouched), and yields precise syntax errors for free. Operator objects
keep \code{evaluate}, \code{d\_da}, and \code{d\_db} and remain
independently testable.
\end{ddbox}

\begin{ddbox}{DD-21}{Incremental public releases}
Each roadmap milestone ships as a public release on PyPI with a Zenodo
DOI, starting with the foundation at \code{v0.1.0}
(Chapter~\ref{ch:roadmap}). The rejected alternative, a single public
release only after every milestone is complete, maximizes time to first
external feedback and concentrates risk for a solo maintainer. What does
not change per release: TDD discipline, the coverage gate, and green CI.
\end{ddbox}

\begin{ddbox}{DD-22}{Solver drivers stay outside \itc{}}
\itc{} is solver-agnostic (NREQ-10). Solver automation packages evolve on
the release cadence of the solvers they drive and carry version-specific
command knowledge that does not belong in a data-management core. Drivers
consume and produce \itc{}-compatible data through \code{itc.load} and the
export formats; embedding any driver was rejected.
\end{ddbox}

\begin{ddbox}{DD-23}{Co-development with pyflightstream}
\itc{} and pyflightstream are developed as consciously integrated sister
libraries: each may generate requirements for the other, and each
documents awareness of the other's architecture. This refines DD-22: the
adapter emitting \itc{}-compatible datasets lives in pyflightstream
behind an optional extra, and \itc{} never imports the driver. Needs the
exporter cannot satisfy with the existing \itc{} surface enter as
candidate requirements carrying a pyflightstream origin. Both
repositories share the role-based review process. Independent evolution
with deferred integration was rejected: requirement flow is cheapest
while both APIs are young.
\end{ddbox}

\begin{ddbox}{DD-24}{Options registry with exact keys}
The central options registry (REQ-104) follows the pandas
\code{register\_option} model with a strict message contract: every
option registers a type, a validator, and a default, and a validator
rejection names the offered value and the accepted domain. Keys are
exact, dot-namespaced strings; partial or abbreviated matching was
rejected because silent prefix resolution can bind a setting to the
wrong option as the registry grows, exactly the quiet misconfiguration
the fail-fast posture (DD-04 family) exists to prevent. The registry
lives in \code{utils/options.py}; the NumPy-only rule is untouched. The
first consumer is the AIAA plot theme as a validated configuration tree
with a frozen testing theme.
\end{ddbox}

\begin{ddbox}{DD-25}{The uncertainties package as a dev-only test oracle}
The \code{uncertainties} package (BSD license) enters the dev dependency
group as an independent oracle for the GUM linear-propagation
mathematics: an oracle test tier cross-validates the random-component
propagation on small analytic cases. It is never a runtime dependency
and is never imported by library code. Adopting it as the propagation
engine was rejected: the two-component model, covariance handling, and
provenance recording require an implementation the library owns, and
the NumPy-only rule bars third-party runtime dependencies in
\code{uncertainty/}.
\end{ddbox}

\begin{ddbox}{DD-26}{scipy as a dev-only geometry test oracle}
The rotation machinery (REQ-38, REQ-101) builds direction-cosine
matrices from the general formulation, composing elementary rotations
in pure NumPy so that \code{core/} keeps the NumPy-only rule and
\itc{} owns the analytical sensitivities \code{dR/dangle} that the
condition-dependent chain-rule propagation needs. \code{scipy} (BSD)
joins the dev dependency group purely as an oracle: a test tier
cross-validates the direction-cosine matrices against
\code{scipy.spatial.transform.Rotation}. It is never a runtime
dependency and is never imported by library code, mirroring DD-25.
Adopting \code{scipy} as the runtime rotation engine was rejected: it
would breach the NumPy-only rule, pull in a heavy runtime dependency,
and still not supply the derivative terms.
\end{ddbox}

\section{Exception Hierarchy}

\begin{table}[H]
\centering
\caption{The \code{ITACAError} hierarchy. Specific classes inherit from
family classes; family classes inherit from \code{ITACAError}.}
\label{tab:exception-hierarchy}
\small
\begin{tabularx}{\textwidth}{@{}llX@{}}
\toprule
\textbf{Family} & \textbf{Specific class} & \textbf{Raised when}\\
\midrule
\code{DataError}        & \code{LoadCoordinateError}        & \code{itc.load} dict-mode tuple length does not match \code{dims}\\
                        & \code{PivotError}                 & \code{db.pivot} called on already-structured VarFrame without \code{dims}\\
                        & \code{PivotDuplicateError}        & \code{db.pivot} finds two datapoints with identical coordinates on every requested dimension\\
                        & \code{DimensionNotFoundError}     & A dimension referenced in an operation is absent\\
                        & \code{VariableNotFoundError}      & A variable referenced in an equation is absent\\
                        & \code{NonNumericDimensionError}   & Numerical operation requested on string-valued dimension\\
                        & \code{SelectionError}             & A coordinate value in \code{select} is not present\\
                        & \code{AxisTranslationError}       & \code{interpolate} axis-translation target is non-monotonic\\
                        & \code{ConcatOverlapError}         & \code{itc.concat} inputs have overlapping coordinates along \code{along}\\
                        & \code{FitDegreeError}            & A moving or global polynomial fit is given too few points for its degree (\code{diff} \code{window <= deg}; \code{smooth} savgol; \code{interpolate} polyfit; \code{fitmodel})\\
                        & \code{GridMismatchError}          & \code{pproc.compare} inputs have different grids and no \code{reference}\\
\addlinespace
\code{ProcessorError}   & \code{ProcessorNotFoundError}     & Unknown processor name passed to \code{itc.pproc}\\
                        & \code{ProcessorValidationError}   & A required variable is missing for the processor\\
                        & \code{ItceqCycleError}            & \code{.itceq} equations contain a cyclic dependency\\
                        & \code{ItceqParseError}            & Generic \code{.itceq} parse failure\\
                        & \code{ProcessorIdempotenceWarning}& Processor reapplied to already-processed data\\
\addlinespace
\code{ProvenanceError}  & \code{DraftModeExportError}       & \code{db.save} on draft VarFrame without \code{allow\_draft=True}\\
                        & \code{OperatingModeMixError}      & Binary op mixes \code{draft} and \code{production} VarFrames\\
                        & \code{HashMismatchError}          & \code{itc.open} detects source-hash drift\\
                        & \code{PipelineCompatibilityError} & \code{Pipeline.apply} target lacks required variables\\
\addlinespace
\code{UncertaintyError} & \code{UncertaintyKeyError}        & \code{set\_uncertainty} key does not match any variable\\
                        & \code{UncertaintyCompatibilityError} & Non-differentiable function in expression with uncertainty inputs\\
                        & \code{CorrelationKeyError}        & \code{set\_correlation} references unknown variable\\
                        & \code{CorrelationMatrixError}     & Provided correlation matrix is not symmetric or violates $|r|\le 1$\\
\addlinespace
\code{AxesError}        & \code{AxisNotFoundError}          & \code{db.rotate} target axis is not registered\\
                        & \code{VectorGroupError}           & \code{db.rotate} cannot resolve a vector group from naming convention\\
                        & \code{RotationMatrixError}        & Provided matrix is not orthogonal\\
\addlinespace
\code{DependencyError}  & \code{MissingDependencyError}     & Optional dependency (matplotlib, plotly, SMT, reportlab) absent\\
\bottomrule
\end{tabularx}
\end{table}
